\chapter*{Abstract}
\addcontentsline{toc}{chapter}{Abstract}

\begin{spacing}{1.5}

% DRAFT: revisar al cerrar los resultados pendientes.
This work presents an on-premise pipeline based on Small Language Models (SLMs) to identify errors in the logs of the Information and Communication Technology Directorate (DTIC) of the University of Cuenca, notify developers and recommend corrective actions. The study uses the \textit{ms-users} microservice corpus (309 files, 27,752,836 events), which a prefilter reduces to 223,104 relevant events grouped into 4,463 batches of 50 events. A batch-level ground truth was built by combining a three-LLM vote (Fleiss' $\kappa=0.886$), review by two experts of 329 batches, rules derived from their criterion and similarity-based propagation, resulting in 130 error batches (2.91\%). The consensus of the three LLMs overestimated errors: only 12.5\% of the batches with a unanimous error vote were errors according to the experts. Eight SLMs were compared with semantic embedding-based evaluation. Against the LLM-consensus reference, granite-4.1-3b obtained the best results (F1 = 0.951, ROC AUC = 0.997); against the expert-corrected ground truth, all models showed high recall (up to 0.984) but low precision, and Qwen3.5-9B achieved the best F1 (0.295). The pipeline integrates Slack notification, a local recommender, error injection and a dashboard; the error-injection experiments remain pending.

\end{spacing}

\hspace{1.27cm}\textit{Author Keywords:} small language models, log analysis, anomaly detection, observability, on-premise, ground truth

\vfill % empuja el bloque CC hacia el pie de página

\begin{minipage}[t]{0.5\linewidth}
  \includegraphics[width=5cm]{figs/logos/image.png}
\end{minipage}
\par\vspace{0.25cm}
\begin{spacing}{0.5}
\begin{footnotesize}\hyphenpenalty=10000 % sin cortar palabras con guion
The content of this work corresponds to the right of expression of the authors and does not compromise the institutional thinking of the University of Cuenca, nor does it release its responsibility before third parties. The \linebreak authors assume responsibility for the intellectual property and copyrights.
\par\vspace{0.6em} % una línea en blanco de separación
\textbf{Institutional Repository:} \href{https://dspace.ucuenca.edu.ec/}{\textcolor{blue}{\underline{\nolinkurl{https://dspace.ucuenca.edu.ec/}}}}
\end{footnotesize}
\end{spacing}
\vspace{0.45cm}\null % pequeño margen entre el enlace y el pie de página
